% \begingroup\expandafter\endgroup\iffalse{}\fi
% This file is not meant to be compiled directly.

\usepackage{atbegshi}% http://ctan.org/pkg/atbegshi
\relax

% \usepackage[sfdefault]{sourcesanspro}
\newcommand{\headingfont}{\sffamily}

\linespread{1.35}  

\usepackage{parskip}
\usepackage[default,osf]{sourcesanspro}

%%%%%%%%%%%%%%%%%%%%%%%%%%%%%%%%%%%%%%%%%%%%%%%%
% Graphics and Tables
% http://en.wikibooks.org/wiki/LaTeX/Importing_Graphics
% http://en.wikibooks.org/wiki/LaTeX/Tables
% http://en.wikibooks.org/wiki/LaTeX/Colors
%%%%%%%%%%%%%%%%%%%%%%%%%%%%%%%%%%%%%%%%%%%%%%%%
% for code blocks
\usepackage{xparse}
\usepackage{xstring}
\usepackage{wrapfig}
\usepackage{tabularx}
\usepackage{setspace}
\usepackage{longtable}
\usepackage{array}
\usepackage{booktabs}
\usepackage{float}
\usepackage[newfloat]{minted}
\setminted{tabsize=2}
\usepackage{listings}
\lstset{frame=tb,
 language=C,
  aboveskip=3mm,
  belowskip=3mm,
  showstringspaces=false,
  columns=flexible,
  basicstyle={\small\ttfamily},
  numbers=none,
  numberstyle=\tiny\color{gray},
  keywordstyle=\color{blue},
  commentstyle=\color{dkgreen},
  stringstyle=\color{mauve},
  breaklines=true,
  breakatwhitespace=true,
  tabsize=3
}

    % load a colour package
\usepackage{xcolor}
\usepackage{refcount}
\usepackage{graphicx}
\usepackage[expansion=false]{microtype}
\usepackage{xurl}
\setlength{\emergencystretch}{1em}
% package to use colour in tables
\usepackage[table]{xcolor}
% Set up how figure and table captions are displayed
\usepackage{caption}
\captionsetup{%
  justification=centering,
  width=0.5\linewidth,
  font=footnotesize,% set font size to footnotesize
  labelfont=bf % bold label (e.g., Figure 3.2) font
}
% Make the standard latex tables look so much better
% Enable the use of frames around, e.g., theorems
% The framed package is used in the example environment
\usepackage{framed}

\usepackage{eurosym}


% Adds support for full page background picture
\usepackage{eso-pic}

%%%%%%%%%%%%%%%%%%%%%%%%%%%%%%%%%%%%%%%%%%%%%%%%
% Mathematics
% http://en.wikibooks.org/wiki/LaTeX/Mathematics
%%%%%%%%%%%%%%%%%%%%%%%%%%%%%%%%%%%%%%%%%%%%%%%%
% Defines new environments such as equation,
% align and split 
\usepackage{amsmath}
% Adds new math symbols
\usepackage{amssymb}
% Use theorems in your document
% The ntheorem package is also used for the example environment
% When using thmmarks, amsmath must be an option as well. Otherwise \eqref doesn't work anymore.
\usepackage[framed,amsmath,thmmarks]{ntheorem}

%%%%%%%%%%%%%%%%%%%%%%%%%%%%%%%%%%%%%%%%%%%%%%%%
% Page Layout
% http://en.wikibooks.org/wiki/LaTeX/Page_Layout
%%%%%%%%%%%%%%%%%%%%%%%%%%%%%%%%%%%%%%%%%%%%%%%%
% Change margins, papersize, etc of the document
\usepackage[
  inner=28mm,% left margin on an odd page
  outer=41mm,% right margin on an odd page
  ]{geometry}
\setlength{\headheight}{14pt}
% Modify how \chapter, \section, etc. look
% The titlesec package is very configureable
\usepackage{titlesec}

\definecolor{accent}{HTML}{211a52}

\setlength{\fboxsep}{4pt}

\newcommand{\secbox}[1]{\colorbox{accent}{\color{white}\hspace{4pt}#1\hspace{4pt}}}

\titleformat{\chapter}[display]
  {\headingfont\huge\bfseries}
  {\colorbox{accent}{\color{white} \thechapter}}
  {1ex}
  {\Huge}

\titleformat{\section}
  {\headingfont\Huge\bfseries}
  {\colorbox{accent}{\color{white} \thesection}}
  {0.7em}
  {}


\titleformat{\subsection}
  {\headingfont\Large\bfseries}
  {\colorbox{accent}{\color{white} \thesubsection}}
  {0.6em}
  {}

\titleformat{\subsubsection}
  {\headingfont\large\bfseries}
  {}
  {0em}
  {}

%\titleformat*{\paragraph}{\normalfont\normalsize\bfseries}
%\titleformat*{\subparagraph}{\normalfont\normalsize\bfseries}

% Clear empty pages between chapters
% \let\origdoublepage\cleardoublepage
% \newcommand{\clearemptydoublepage}{%
%   \clearpage
%   {\pagestyle{empty}\origdoublepage}%
% }
% \let\cleardoublepage\clearemptydoublepage

% Change the headers and footers
\usepackage{fancyhdr}
\usepackage{fancyvrb}
\pagestyle{fancy}
\fancyhf{} %delete everything
\renewcommand{\headrulewidth}{0pt} %remove the horizontal line in the header
\fancyhead[R]{\small\nouppercase\leftmark} %even page - chapter title
\fancyhead[LO]{\small\nouppercase\rightmark} %uneven page - section title
\fancyhead[L,RO]{\thepage} %page number on all pages
% Do not stretch the content of a page. Instead,
% insert white space at the bottom of the page
\raggedbottom
% Enable arithmetics with length. Useful when
% typesetting the layout.
\usepackage{calc}

\usepackage{pageslts}

%%%%%%%%%%%%%%%%%%%%%%%%%%%%%%%%%%%%%%%%%%%%%%%%
% Bibliography
% http://en.wikibooks.org/wiki/LaTeX/Bibliography_Management
%%%%%%%%%%%%%%%%%%%%%%%%%%%%%%%%%%%%%%%%%%%%%%%%
\usepackage[
  backend=biber,
  bibencoding=utf8,
  style=numeric-comp,
  sorting=none,
  defernumbers=true
]{biblatex}
\addbibresource{report/bib/mybib.bib}

%%%%%%%%%%%%%%%%%%%%%%%%%%%%%%%%%%%%%%%%%%%%%%%%
% Misc
%%%%%%%%%%%%%%%%%%%%%%%%%%%%%%%%%%%%%%%%%%%%%%%%
% Add bibliography and index to the table of
% contents
\usepackage[nottoc]{tocbibind}
% Add the command \pageref{LastPage} which refers to the
% page number of the last page
% \usepackage{lastpage}
% Add todo notes in the margin of the document
\setlength{\marginparwidth}{2cm}
\usepackage[
%  disable, %turn off todonotes
  colorinlistoftodos, %enable a coloured square in the list of todos
  textwidth=\marginparwidth, %set the width of the todonotes
  textsize=scriptsize, %size of the text in the todonotes
  ]{todonotes}

%%%%%%%%%%%%%%%%%%%%%%%%%%%%%%%%%%%%%%%%%%%%%%%%
% Hyperlinks
% http://en.wikibooks.org/wiki/LaTeX/Hyperlinks
%%%%%%%%%%%%%%%%%%%%%%%%%%%%%%%%%%%%%%%%%%%%%%%%
% Enable hyperlinks and insert info into the pdf
% file. Hypperref should be loaded as one of the 
% last packages
\usepackage{hyperref}
\hypersetup{%
    hypertexnames=false,%
    pdfinfo={%
    Title={\reporttitle},%
    Author={\authors},%
    Subject={\subject},%
    Producer={pdfLaTeX},%
    Creator={\authors},%
    % Keywords={},%    
    }%   
}
\usepackage{pdfpages}
\usepackage{subcaption}

\usepackage{letltxmacro}
\LetLtxMacro{\oldautoref}{\autoref}
\renewcommand{\autoref}[1]{\textbf{\oldautoref{#1}}}

\usepackage{blindtext}


\usepackage{lastpage}
\usepackage{totcount}
\regtotcounter{page}

\usepackage{copyrightbox}

\usepackage{multicol}

% Source Code float
\floatstyle{plaintop}
\newfloat{source}{H}{los}[section]
\floatname{source}{Source Code}
\renewcommand{\thesource}{\arabic{source}}
\newenvironment{sourceblock}{\captionsetup{type=source}}{}
\SetupFloatingEnvironment{source}{name=Source Code}

% Code Block float
\floatstyle{plaintop}
\newfloat{code}{H}{loc}[section]
\floatname{code}{Code Block}
\renewcommand{\thesource}{\arabic{source}}
\newenvironment{codeblock}{\captionsetup{type=code}}{}
\SetupFloatingEnvironment{code}{name=Code Block}


